\section*{Problem 389. A divisibility problem for two consecutive blocks}

\subsection*{1) ROUND-2 OBJECTIVE}

Path (A) was not realistic after Round~1, because no general existence theorem for all $n$ had been found.  The most promising Round-2 goal is therefore the fail-safe route: prove theorems that strictly advance beyond Round~1.

I will do this in three ways.

\begin{itemize}
\item I derive an explicit digit-sum criterion from the Round-1 valuation reformulation.
\item I give infinite families of solutions for $n=2$ and $n=3$ (and the trivial complete solution for $n=1$).
\item I prove that for every fixed $n$, infinitely many $k$ fail.
\end{itemize}

So the global problem remains open, but the picture is now much sharper.

\subsection*{2) Round-1 FOUNDATION USED}

I use the following Round-1 facts.

\begin{itemize}
\item The problem asks whether for every $n\ge 1$ there exists $k\ge 1$ such that
\[
\prod_{i=0}^{k-1}(n+i)\mid \prod_{i=0}^{k-1}(n+k+i).
\]

\item Round-1 binomial reformulation:
\[
\prod_{i=0}^{k-1}(n+i)\mid \prod_{i=0}^{k-1}(n+k+i)
\quad\Leftrightarrow\quad
{n+k-1\choose k}\mid {n+2k-1\choose k}.
\]

\item Round-1 valuation reformulation:
the divisibility is equivalent to
\[
v_p\!\left((n+2k-1)!\right)+v_p\!\left((n-1)!\right)\ge 2v_p\!\left((n+k-1)!\right)
\]
for every prime $p$.
\end{itemize}

\subsection*{3) NEW INSIGHT / TOOL (ROUND-2)}

The key new tool is a \emph{digit-sum criterion} obtained from the Round-1 valuation identity via Legendre's formula.

If $s_p(m)$ denotes the sum of the base-$p$ digits of $m$, then
\[
v_p(t!)=\frac{t-s_p(t)}{p-1}.
\]
Substituting this into the Round-1 valuation condition yields an exact criterion in terms of
\[
2s_p(n+k-1)-s_p(n+2k-1)-s_p(n-1).
\]

This criterion is strong enough to prove:
\begin{itemize}
\item every $k$ works for $n=1$;
\item infinitely many $k$ work for $n=2$;
\item infinitely many $k$ work for $n=3$;
\item for each fixed $n$, infinitely many $k$ do \emph{not} work.
\end{itemize}

\subsection*{4) ATTACK PLAN (ROUND-2)}

The exact Round-1 gap was the absence of any constructive mechanism producing or ruling out admissible $k$.

To overcome that, I first derive the digit-sum test.  Then I use it in two opposite directions:

\begin{enumerate}
\item choose $k$ so that the relevant base-$p$ digit sums force divisibility (positive families);
\item choose $k$ so that one particular prime gives a negative digit-sum expression (negative families).
\end{enumerate}

This supplies strict, rigorous progress on both the existence and nonexistence sides.

\subsection*{5) WORK (ROUND-2)}

Define
\[
T_n:=\left\{k\ge 1:\ \prod_{i=0}^{k-1}(n+i)\mid \prod_{i=0}^{k-1}(n+k+i)\right\}.
\]

\medskip
\noindent\textbf{Claim 1 (digit-sum criterion).}
\emph{A positive integer $k$ belongs to $T_n$ if and only if for every prime $p$ dividing $\prod_{i=0}^{k-1}(n+i)$ one has}
\[
2s_p(n+k-1)-s_p(n+2k-1)-s_p(n-1)\ge 0.
\]

\medskip
\noindent\textbf{Proof.}
By Round-1, $k\in T_n$ if and only if
\[
v_p\!\left((n+2k-1)!\right)+v_p\!\left((n-1)!\right)-2v_p\!\left((n+k-1)!\right)\ge 0
\]
for every prime $p$.

Now apply Legendre's formula
\[
v_p(t!)=\frac{t-s_p(t)}{p-1}.
\]
The left-hand side becomes
\[
\frac{(n+2k-1)-s_p(n+2k-1)+(n-1)-s_p(n-1)-2(n+k-1)+2s_p(n+k-1)}{p-1},
\]
which simplifies to
\[
\frac{2s_p(n+k-1)-s_p(n+2k-1)-s_p(n-1)}{p-1}.
\]
Therefore the valuation inequality is equivalent to
\[
2s_p(n+k-1)-s_p(n+2k-1)-s_p(n-1)\ge 0.
\]

It is enough to check primes dividing the left block $\prod_{i=0}^{k-1}(n+i)$, because if a prime $p$ does not divide that block, then its left-side valuation is $0$ and the divisibility condition at $p$ is automatic.
\hfill\textit{QED.}

\medskip
\noindent\textbf{Claim 2.}
\emph{$T_1=\{1,2,3,\dots\}$.  In other words, every $k\ge 1$ works for $n=1$.}

\medskip
\noindent\textbf{Proof.}
For $n=1$ the divisibility becomes
\[
1\cdot 2\cdots k \mid (k+1)(k+2)\cdots(2k).
\]
The left-hand side is $k!$, and the right-hand side is
\[
\frac{(2k)!}{k!}={2k\choose k}\,k!.
\]
Hence the right-hand side is an integer multiple of $k!$, so every $k$ works.
\hfill\textit{QED.}

\medskip
\noindent\textbf{Claim 3.}
\emph{$T_2$ is infinite.  More precisely, if $p>5$ and $q$ are primes with}
\[
\frac{3p}{2}<q<2p,
\]
\emph{then}
\[
k=pq-1
\]
\emph{belongs to $T_2$.}

\medskip
\noindent\textbf{Proof.}
For $n=2$, Round-1 gives
\[
k\in T_2
\quad\Leftrightarrow\quad
{k+1\choose k}=k+1\mid {2k+1\choose k}.
\]
So it suffices to show that $k+1=pq$ divides ${2k+1\choose k}$.

Because $q<2p$, the prime factorization of $k+1$ is exactly $pq$ with two distinct primes $p$ and $q$.
Hence it is enough to show that both $p$ and $q$ divide ${2k+1\choose k}$.

By Claim~1 with $n=2$, it is enough to verify
\[
2s_r(k+1)-s_r(2k+1)-1\ge 0
\]
for $r=p$ and $r=q$.

\smallskip
\noindent\emph{Base $p$.}
Since $k+1=pq$ and $q<2p$, we can write
\[
k+1=pq=p^2+(q-p)p.
\]
Therefore
\[
s_p(k+1)=1+(q-p).
\]
Also
\[
2k+1=2pq-1=3p^2+(2q-3p-1)p+(p-1).
\]
The coefficient $2q-3p-1$ is a valid base-$p$ digit because $q>3p/2$ implies $2q-3p\ge 1$, and $q<2p$ implies $2q-3p-1\le p-2$.
Hence
\[
s_p(2k+1)=3+(2q-3p-1)+(p-1)=2q-2p+1.
\]
Thus
\[
2s_p(k+1)-s_p(2k+1)-1
=2(q-p+1)-(2q-2p+1)-1=0.
\]

\smallskip
\noindent\emph{Base $q$.}
Now
\[
k+1=pq=p\cdot q+0,
\]
so
\[
s_q(k+1)=p.
\]
Also
\[
2k+1=2pq-1=q^2+(2p-q-1)q+(q-1).
\]
Again the coefficient $2p-q-1$ is a valid base-$q$ digit because $q<2p$ implies $2p-q-1\ge 0$, and $q>p$ implies $2p-q-1<q$.
Therefore
\[
s_q(2k+1)=1+(2p-q-1)+(q-1)=2p-1,
\]
and hence
\[
2s_q(k+1)-s_q(2k+1)-1=2p-(2p-1)-1=0.
\]

So both primes $p$ and $q$ divide ${2k+1\choose k}$, hence $pq=k+1$ divides ${2k+1\choose k}$, and therefore $k\in T_2$.

Finally, there are infinitely many such prime pairs $(p,q)$.  Indeed, Nagura's theorem gives a prime in $(x,6x/5)$ for every sufficiently large $x$; taking $x=3p/2$ yields a prime
\[
q\in\left(\frac{3p}{2},\frac{9p}{5}\right)\subset\left(\frac{3p}{2},2p\right)
\]
for every sufficiently large prime $p$.
So $T_2$ is infinite.
\hfill\textit{QED.}

\medskip
\noindent\textbf{Claim 4.}
\emph{$T_3$ is infinite.  More precisely, if $k\in T_2$, then $k-1\in T_3$.}

\medskip
\noindent\textbf{Proof.}
Assume $k\in T_2$.
Then by the $n=2$ binomial reformulation,
\[
k+1\mid {2k+1\choose k}.
\]
Use the standard identity
\[
{2k+1\choose k}=\frac{2k+1}{k+1}{2k\choose k}.
\]
Multiplying the divisibility $k+1\mid {2k+1\choose k}$ by $k$ gives
\[
k(k+1)\mid k{2k+1\choose k}=(2k+1){2k\choose k-1}.
\]
Now
\[
\gcd(2k+1,k)=1,\qquad \gcd(2k+1,k+1)=1,
\]
so
\[
\gcd(2k+1,k(k+1))=1.
\]
Therefore
\[
k(k+1)\mid {2k\choose k-1}.
\]
Since $k(k+1)$ is always even, this implies
\[
{k+1\choose 2}=\frac{k(k+1)}{2}\mid {2k\choose k-1}.
\]
Now apply the Round-1 binomial reformulation with $n=3$ and $K:=k-1$:
\[
K\in T_3
\quad\Leftrightarrow\quad
{K+2\choose K}={K+2\choose 2}\mid {2K+2\choose K}.
\]
Substituting $K=k-1$ gives exactly
\[
{k+1\choose 2}\mid {2k\choose k-1},
\]
which we have proved.  Hence $k-1\in T_3$.
\hfill\textit{QED.}

Combining Claims~3 and~4, the explicit family
\[
k=pq-2
\]
belongs to $T_3$ whenever $p>5$ and $q$ are primes with $3p/2<q<2p$.

\medskip
\noindent\textbf{Claim 5.}
\emph{For every fixed $n\ge 2$, the complement of $T_n$ is infinite.  In fact, if $p>n-1$ is prime and $m\ge 2$, then}
\[
k=2p^m-n+1
\]
\emph{does not belong to $T_n$.}

\medskip
\noindent\textbf{Proof.}
Fix $n\ge 2$, choose a prime $p>n-1$, choose $m\ge 2$, and define
\[
k=2p^m-n+1.
\]
Then
\[
n+k-1=2p^m,
\]
so $p$ divides the left block $\prod_{i=0}^{k-1}(n+i)$, and Claim~1 applies to this prime $p$.

Because $p>n-1$, the base-$p$ expansion of $n-1$ is a single digit, so
\[
s_p(n-1)=n-1.
\]
Also
\[
n+k-1=2p^m,
\]
hence
\[
s_p(n+k-1)=2.
\]

Now
\[
n+2k-1=4p^m-n+1.
\]
Write this in base $p$ as
\[
4p^m-n+1
=3p^m+(p-1)(p+p^2+\cdots+p^{m-1})+(p-n+1).
\]
This is valid because $p>n-1$, so $0\le p-n+1\le p-1$.
Therefore
\[
s_p(n+2k-1)=3+(m-1)(p-1)+(p-n+1).
\]

Insert these values into Claim~1:
\[
2s_p(n+k-1)-s_p(n+2k-1)-s_p(n-1)
\]
\[
=4-\bigl(3+(m-1)(p-1)+(p-n+1)\bigr)-(n-1)
=1-p-(m-1)(p-1)<0.
\]
So the digit-sum criterion fails for the prime $p$, and therefore $k\notin T_n$.

Since there are infinitely many choices of $p$ and $m$, the complement of $T_n$ is infinite.
\hfill\textit{QED.}

\subsection*{6) ADVERSARIAL VERIFICATION}

I check the new arguments against the main possible flaws.

\begin{itemize}
\item \textbf{Index shift.}
The literature on this problem often uses the parameter $n-1$ instead of $n$.
To avoid any hidden mismatch, every statement above was re-derived in the present notation.

\item \textbf{Claim 1 uses only primes dividing the left block.}
That is correct: divisibility needs to be checked only at primes that occur on the left.  For primes not dividing the left block, the left valuation is $0$ and the divisibility condition is automatic.

\item \textbf{Base-$p$ digit expansions in Claim~3.}
The inequalities
\[
\frac{3p}{2}<q<2p
\]
were used exactly to ensure that the coefficients in the base-$p$ and base-$q$ expansions are valid digits.  I checked both bounds explicitly in the proof.

\item \textbf{Infinitely many prime pairs $(p,q)$.}
This does not follow from Bertrand alone, but it does follow from Nagura's theorem on primes in short intervals, which is why I used Nagura explicitly.

\item \textbf{The reduction from $T_2$ to $T_3$.}
The step
\[
k(k+1)\mid (2k+1){2k\choose k-1}
\quad\Rightarrow\quad
k(k+1)\mid {2k\choose k-1}
\]
uses $\gcd(2k+1,k(k+1))=1$, and this was checked directly.

\item \textbf{Claim~5.}
The chosen prime $p$ really does divide the left block because $n+k-1=2p^m$ is the last term of that block.  The base-$p$ expansion of $4p^m-n+1$ was checked by direct summation.
\end{itemize}

\subsection*{7) FINAL}

\textbf{UNRESOLVED (BUT STRICTLY ADVANCED).}

The strict advances beyond Round~1 are the following.

\begin{enumerate}
\item There is an exact digit-sum criterion for membership in $T_n$.

\item The problem is completely solved for $n=1$: every $k$ works, so in this case the complement of $T_1$ is empty.

\item There are infinitely many solutions for $n=2$, for example
\[
k=pq-1
\]
with primes $p>5$ and $3p/2<q<2p$.

\item There are infinitely many solutions for $n=3$, for example
\[
k=pq-2
\]
with the same prime-pair condition.

\item For every fixed $n\ge 2$, there are infinitely many failures; indeed
\[
k=2p^m-n+1
\]
never works whenever $p>n-1$ is prime and $m\ge 2$.
\end{enumerate}

What remains open is the original global existence question:
for each $n\ge 1$, must the set $T_n$ be nonempty?

\subsection*{8) COMPLETION ESTIMATE (MANDATORY)}

COMPLETION: 45\%

\subsection*{9) REFERENCES}

\begin{enumerate}
\item M. Ulas, \emph{On the Erd\H{o}s--Straus and Erd\H{o}s--Graham problems}, especially the digit-sum criterion and the explicit families for the shifted indexing version of this problem.

\item J. Nagura, \emph{On the interval containing at least one prime number}, for the prime-in-short-interval statement used to make the family in Claim~3 infinite.

\item The Erd\H{o}s Problems site, Problem~389, and OEIS A375071 for computational background only.
\end{enumerate}



